\chapter{Kafka in the Cluster: Strimzi and the Mail Flow}
\label{ch:strimzi}

Milestone~6 moves the broker of Chapter~10 from Compose into k3s and
deploys the last service --- notification-service --- so the first
end-to-end \emph{event} flow runs in the cluster: register a user,
receive a mail. It is also the chapter where Chapter~14's reconciliation
idea grows its most important extension: the \emph{operator}.

\section{Operators and custom resources}

Kubernetes ships controllers for its built-in kinds --- Deployments,
Services. An \emph{operator} is a controller for kinds you add: a
\emph{CustomResourceDefinition} (CRD) teaches the API server a new
noun (\code{Kafka}, \code{KafkaTopic}), and the operator's
reconciliation loop turns instances of it into ordinary objects ---
pods, Services, ConfigMaps --- and keeps them converged. Operational
knowledge that used to live in runbooks (how to roll a broker safely,
how to compute advertised listeners, how to add a partition) is
encoded in the operator's code. Strimzi is the Kafka operator; the
pattern is universal (Postgres operators, Prometheus operator, cert-manager).

\section{Installing the operator: an admin act}

\begin{lstlisting}[language=bash]
# deploy/k8s/install-strimzi.sh (abridged)
kubectl apply --server-side --force-conflicts \
  -f "https://strimzi.io/install/latest?namespace=${NS}" -n "$NS"   (*@\co{1}@*)
kubectl -n "$NS" rollout status deployment/strimzi-cluster-operator
\end{lstlisting}

\begin{description}
  \item[\co{1}] The bundle contains the CRDs, the operator Deployment,
  and its RBAC --- much of it \emph{cluster-scoped}. Two consequences:
  it is applied by an admin, not by Jenkins (whose Role, Chapter~22,
  stops at the namespace), and it is applied \emph{server-side}
  because the CRD documents are too large for the client-side
  last-applied annotation. Idempotent, so it doubles as the upgrade
  command. The \code{?namespace=} parameter scopes the operator to
  watch \code{ts} only.
\end{description}

The kustomization (Chapter~20) lists \code{kafka.yaml}, but that file
cannot apply until this script has run once --- the third member of the
``admin prerequisites'' set, after the namespace and the secrets.

\begin{pitfall}[no matches for kind ``Kafka'' in version v1beta2]
Symptom: applying a Kafka manifest copied from any tutorial fails with
this error even though the operator is installed and healthy. Cause:
Strimzi~1.x serves only \code{kafka.strimzi.io/v1}; the
\code{v1beta2} version that a decade of blog posts use is gone.
Check what a CRD actually serves before trusting an example:
\code{kubectl get crd kafkas.kafka.strimzi.io -o jsonpath=\{.spec.versions[*].name\}}.
This project hit it live on installation day.
\end{pitfall}

\section{The cluster, declared}

\begin{lstlisting}[language=yaml]
apiVersion: kafka.strimzi.io/v1
kind: KafkaNodePool                        (*@\co{1}@*)
metadata:
  name: dual-role
  labels: { strimzi.io/cluster: ts-kafka } (*@\co{2}@*)
spec:
  replicas: 1
  roles: [controller, broker]
  storage: { type: persistent-claim, size: 5Gi, deleteClaim: false } (*@\co{3}@*)
  resources:
    requests: { memory: 768Mi, cpu: 200m }
    limits: { memory: 1Gi }
---
apiVersion: kafka.strimzi.io/v1
kind: Kafka
metadata:
  name: ts-kafka
spec:
  kafka:
    listeners:
      - { name: plain, port: 9092, type: internal, tls: false }  (*@\co{4}@*)
    config:
      offsets.topic.replication.factor: 1  (*@\co{5}@*)
      default.replication.factor: 1
      min.insync.replicas: 1
      log.retention.hours: 168
    jvmOptions: { -Xms: 256m, -Xmx: 512m } (*@\co{6}@*)
  entityOperator:
    topicOperator: {}                      (*@\co{7}@*)
\end{lstlisting}

\begin{description}
  \item[\co{1}] Two objects: the \emph{pool} says where Kafka runs
  (replicas, roles, storage, resources); the \code{Kafka} resource says
  how it behaves. One dual-role KRaft node is the Kubernetes spelling of
  Compose's \code{process.roles=broker,controller}.
  \item[\co{2}] The label joins pool to cluster; the operator matches on
  it. A mismatch is silent --- the pool simply never materializes.
  \item[\co{3}] A PVC via local-path, exactly like the databases
  (Chapter~19). \code{deleteClaim: false}: deleting the Kafka resource
  keeps the log --- the same asymmetry StatefulSets have by default.
  \item[\co{4}] Compare with the three-listener, two-advertised-address
  block in \code{docker-compose.yml}: gone. You state \emph{one
  internal listener}; the operator computes bootstrap and per-broker
  Services, the advertised addresses, and the KRaft controller listener.
  Chapter~10's maze was operational knowledge; it now lives in the
  operator.
  \item[\co{5}] Single broker means every replication factor must be 1,
  or the first topic creation fails asking for brokers that do not
  exist.
  \item[\co{6}] Kafka's default heap assumes a server; half a GiB
  serves a homelab trickle and leaves memory for the JVM services.
  \item[\co{7}] The entity operator's topic half watches
  \code{KafkaTopic} objects and reconciles them into real topics. The
  user half is omitted: no authentication, no users to manage.
\end{description}

Strimzi names the bootstrap Service \code{<cluster>-kafka-bootstrap};
every producer and consumer therefore gets
\code{SPRING\_KAFKA\_BOOTSTRAP\_SERVERS=ts-kafka-kafka-bootstrap:9092}
in its Deployment --- the \code{localhost:9092} config default of
Part~III overridden the way every other address was.

\begin{decision}[operator, not a hand-written StatefulSet]
A plain KRaft StatefulSet (one pod, one PVC, the compose env block
copied in) would cost less memory --- no operator, no entity operator
--- and it is the honest fallback if CPU contention ever bites. The
operator wins because it is what Kubernetes shops run (Red Hat sells it
as AMQ Streams), because topics become declarative objects, and because
broker upgrades and rolling restarts become someone else's tested code.
Learning it is the point; the fallback stays documented in
\code{STACK-K3S-JENKINS.md}.
\end{decision}

\section{Topics as objects}

\begin{lstlisting}[language=yaml]
apiVersion: kafka.strimzi.io/v1
kind: KafkaTopic
metadata:
  name: ts.user.registered
  labels: { strimzi.io/cluster: ts-kafka }
spec: { partitions: 1, replicas: 1 }
\end{lstlisting}

Kafka would auto-create these on first use; declaring them makes
partition count a reviewed decision, removes the day-one race between
producer and topic, and turns \code{kubectl get kafkatopics} into the
inventory of every event contract in the system (four, matching
Chapter~6's table). One partition per topic: total order, which
trivially preserves the per-user ordering the publishers key for.
Partitions can be added later, never removed.

\section{The last service, and its mailbox}

notification-service's manifest is the standard Deployment of
Chapter~15 with three specifics: the bootstrap address above;
\code{SPRING\_MAIL\_HOST=maildev} / \code{PORT=1025} overriding the
config's \code{\$\{MAIL\_HOST:localhost\}} default via relaxed binding;
and probes on \code{/actuator/health} --- this service has no Spring
Security at all, so the endpoint is open by default (the third answer
to Chapter~17's probe question). Maildev is a tiny Deployment plus
Service (ports 1025 SMTP, 1080 web); it delivers nothing and shows
everything it caught. No Ingress: it is an operator's tool, reached by
\code{port-forward}.

\begin{pitfall}[NotReady --- FatalProblem, while the pod is Running]
Symptom: minutes after applying, \code{kubectl get kafka} shows
\code{READY} empty and a condition reading ``Pod is unschedulable or is
not starting'' --- yet \code{get pods} shows the broker \code{1/1
Running}. Cause: the operator's reconciliation timed out while the
600\,MB Kafka image was still pulling over Wi-Fi, and wrote a failure
status; the \emph{next} periodic loop (every two minutes) observed the
running pod and flipped it to Ready. Operator status is eventually
honest. Before reacting to a red condition, check the pod it is
describing --- and check the operator log, where
``reconciled'' lines show the loops still turning.
\end{pitfall}

\begin{hardware}
Operator $+$ broker $+$ entity operator cost about 1.5\,GB here, with
the broker's heap capped at 512\,MiB. The first image pull is the
expensive part --- expect ten minutes on residential Wi-Fi, and the
stale-status pitfall above. Everything afterwards is fast: the image
is cached on the node, which the smoke test below exploits.
\end{hardware}

\begin{handson}[prove the broker, then prove the flow]
A throwaway pod with the Kafka CLI, in the same network position as
the services, using the image the node already holds:
\begin{lstlisting}[language=bash]
kubectl -n ts run kafka-smoke --rm -i --restart=Never \
  --image=quay.io/strimzi/kafka:1.2.0-kafka-4.3.1 -- sh -c '
  B=ts-kafka-kafka-bootstrap:9092
  bin/kafka-topics.sh --bootstrap-server $B --create --topic smoke \
      --partitions 1 --replication-factor 1
  echo hello | bin/kafka-console-producer.sh --bootstrap-server $B --topic smoke
  bin/kafka-console-consumer.sh --bootstrap-server $B --topic smoke \
      --from-beginning --max-messages 1 --timeout-ms 20000
  bin/kafka-topics.sh --bootstrap-server $B --delete --topic smoke'
\end{lstlisting}
Use a scratch topic, not a real one: a plain-text ``hello'' on
\code{ts.user.registered} is a poison pill for the JSON consumer
(Chapter~11) --- skipped gracefully thanks to
\code{ErrorHandlingDeserializer}, but why find out in production. Then
the real flow, once the pipeline has deployed the services:
\begin{lstlisting}[language=bash]
kubectl -n ts port-forward svc/maildev 1080:1080 &
curl -X POST http://ts.local/api/auth/register -H 'Content-Type: application/json' \
  -d '{"email":"t@example.com","password":"Passw0rd!","fullName":"Test"}'
# open http://localhost:1080 -- the activation mail is there
kubectl -n ts logs deploy/notification-service | grep -i registered
\end{lstlisting}
That request crosses every part of this book: Ingress $\rightarrow$
gateway $\rightarrow$ auth-service $\rightarrow$ Postgres $\rightarrow$
Kafka $\rightarrow$ notification-service $\rightarrow$ SMTP.
\end{handson}

\section{Outside the project}

Managed Kafka --- AWS MSK, Confluent Cloud, Azure Event Hubs' Kafka
endpoint --- is the operator's operator: the same \code{Kafka}-shaped
declaration made in a cloud console, with brokers you never see.
Migration from here is a bootstrap string plus SASL/TLS on the
listener; the \code{KafkaTopic} objects would give way to Terraform
resources saying exactly the same thing. The other lesson to carry: the
operator pattern is how Kubernetes hosts \emph{any} stateful platform
component --- Postgres (CloudNativePG), Redis, Elasticsearch --- and
having installed, mis-versioned, and status-debugged one, you know the
shape of all of them.

\begin{quiz}
\begin{enumerate}
  \item Define an operator in terms of Chapter~14's controllers, and
  name the two kinds of object Strimzi added to this cluster.
  \item Why is \code{install-strimzi.sh} an admin step while
  \code{kafka.yaml} is applied by Jenkins? Which RBAC concept draws the
  line?
  \item List three things Compose's Kafka block configured by hand that
  the \code{Kafka} resource does not mention, and say who computes them
  now.
  \item You see \code{NotReady / FatalProblem} on the Kafka resource.
  What two things do you check before acting, and why might the status
  be wrong?
  \item Why declare topics as objects when Kafka auto-creates them, and
  why one partition each?
\end{enumerate}
\end{quiz}
